
\documentclass{article}
\usepackage{colortbl}
\usepackage{makecell}
\usepackage{multirow}
\usepackage{supertabular}

\begin{document}

\newcounter{utterance}

\centering \large Interaction Transcript for game `wordle', experiment `medium\_frequency\_words\_no\_clue\_no\_critic', episode 3 with NEW\_Reflect\_Agent.
\vspace{24pt}

{ \footnotesize  \setcounter{utterance}{1}
\setlength{\tabcolsep}{0pt}
\begin{supertabular}{c@{$\;$}|p{.15\linewidth}@{}p{.15\linewidth}p{.15\linewidth}p{.15\linewidth}p{.15\linewidth}p{.15\linewidth}}
    \# & \multicolumn{2}{c}{Player} && \multicolumn{2}{c}{Game Master} \\
    \hline

    \theutterance \stepcounter{utterance}  
    & & & \multicolumn{4}{p{0.6\linewidth}}{
        \cellcolor[rgb]{0.9,0.9,0.9}{
            \makecell[{{p{\linewidth}}}]{
                \texttt{\tiny{[P1$\langle$GM]}}
                \texttt{You are a language wizard who likes to guess words by using the given rules.} \\
\\ 
\texttt{Welcome to Wordle! You have six attempts to guess the target word, a valid English word of five lowercase letters (a{-}z). Please use the tags "explanation:" and "guess:" to provide a concise explanation for each guess.} \\
\\ 
\texttt{For instance, if your guess is "apple", your response should be} \\
\texttt{explanation: this is a common five{-}letter English word, and I am starting my guess with this word.} \\
\texttt{guess: apple} \\
\\ 
\texttt{After each guess, your answer will be validated, and you will receive feedback indicating which letters are correct (green), which letters are correct but in the wrong position (yellow), and which letters are incorrect (red). This feedback can be useful in determining which letters to include or exclude in your next guess.} \\
\\ 
\texttt{For example, the feedback for "apple" might be:} \\
\texttt{guess\_feedback: a<yellow> p<yellow> p<green> l<yellow> e<red>} \\
\\ 
\texttt{The explanation should contain details about how the guess\_feedback is used to arrive at a new guess.} \\
\\ 
\texttt{Do not generate any other text.} \\
\\ 
\texttt{Let's begin with your first guess.} \\
            }
        }
    }
     \\ \\

    \theutterance \stepcounter{utterance}  
    & \multicolumn{4}{p{0.45\linewidth}}{
        \cellcolor[rgb]{0.95,0.95,0.95}{
            \makecell[{{p{\linewidth}}}]{
                \texttt{\tiny{[P1$|internal]}}
                \texttt{Reflections from previous episodes:} \\
\texttt{{-}} \\
\texttt{ Failed. I repeated the same guess after identical feedback, wasting an attempt. The feedback showed 'a' and 'r' were yellow but misplaced, yet I didn't systematically test new positions or unused common vowels. Next time, I would use the feedback to make a distinct guess that repositions yellow letters and tests new consonants.} \\
\texttt{{-}} \\
\texttt{ I succeeded in guessing the target word within six attempts. My initial guess "stare" efficiently tested common vowels and consonants, and I used the feedback systematically to narrow down possibilities by prioritizing green letters, repositioning yellows, and eliminating incorrect letters. I maintained a logical progression without repeating excluded letters. Next time, I would follow the same approach—starting with a strong word like "stare" and adapting precisely to the feedback.} \\
\texttt{{-}} \\
\texttt{ Failed. I repeated the same guess "crane" multiple times despite identical feedback, wasting attempts. I should have used the feedback to adjust my guess, especially noting that 'a' is yellow (correct letter, wrong position) and 'c' is yellow, while other letters are red. Next time, I will vary guesses systematically to test new letter positions and exclude red letters.} \\
            }
        }
    }
    & & \\ \\

    \theutterance \stepcounter{utterance}  
    & \multicolumn{4}{p{0.45\linewidth}}{
        \cellcolor[rgb]{0.95,0.95,0.95}{
            \makecell[{{p{\linewidth}}}]{
                \texttt{\tiny{[P1$|internal]}}
                \texttt{{[}attempt 1 – valid{]}} \\
\texttt{explanation: I start with "stare" because it contains common vowels (a, e) and frequently used consonants (s, t, r), allowing me to efficiently gather feedback on letter presence and position.} \\
\texttt{guess: stare} \\
            }
        }
    }
    & & \\ \\

    \theutterance \stepcounter{utterance}  
    & \multicolumn{4}{p{0.6\linewidth}}{
        \cellcolor[rgb]{0.561,0.737,0.561}{
            \makecell[{{p{\linewidth}}}]{
                \texttt{\tiny{[P1$\rangle$GM]}}
                \texttt{explanation: I start with "stare" because it contains common vowels (a, e) and frequently used consonants (s, t, r), allowing me to efficiently gather feedback on letter presence and position.} \\
\texttt{guess: stare} \\
            }
        }
    }
    & & \\ \\

    \theutterance \stepcounter{utterance}  
    & & & \multicolumn{2}{p{0.3\linewidth}}{
        \cellcolor[rgb]{0.9,0.9,0.9}{
            \makecell[{{p{\linewidth}}}]{
                \texttt{\tiny{[GM$|$GM]}}
                \texttt{attempts = 1} \\
\texttt{target\_word = crank} \\
\texttt{guess = stare} \\
\texttt{guess\_feedback = s<red> t<red> a<green> r<yellow> e<red>} \\
            }
        }
    }
    & & \\ \\

    \theutterance \stepcounter{utterance}  
    & & & \multicolumn{4}{p{0.6\linewidth}}{
        \cellcolor[rgb]{0.9,0.9,0.9}{
            \makecell[{{p{\linewidth}}}]{
                \texttt{\tiny{[P1$\langle$GM]}}
                \texttt{guess\_feedback: s<red> t<red> a<green> r<yellow> e<red>} \\
\\ 
\texttt{Provide your response only in this format.} \\
\texttt{explanation: details} \\
\texttt{guess: word} \\
            }
        }
    }
     \\ \\

    \theutterance \stepcounter{utterance}  
    & \multicolumn{4}{p{0.45\linewidth}}{
        \cellcolor[rgb]{0.95,0.95,0.95}{
            \makecell[{{p{\linewidth}}}]{
                \texttt{\tiny{[P1$|internal]}}
                \texttt{{[}attempt 1 – valid{]}} \\
\texttt{explanation: The feedback shows that 'a' is in the correct position (green), 'r' is present but in the wrong position (yellow), and 's', 't', 'e' are incorrect (red). I will keep 'a' in the first position, include 'r' in the next guess but not in the same position, and avoid 's', 't', 'e'.  } \\
\texttt{guess: rater} \\
            }
        }
    }
    & & \\ \\

    \theutterance \stepcounter{utterance}  
    & \multicolumn{4}{p{0.6\linewidth}}{
        \cellcolor[rgb]{0.561,0.737,0.561}{
            \makecell[{{p{\linewidth}}}]{
                \texttt{\tiny{[P1$\rangle$GM]}}
                \texttt{explanation: The feedback shows that 'a' is in the correct position (green), 'r' is present but in the wrong position (yellow), and 's', 't', 'e' are incorrect (red). I will keep 'a' in the first position, include 'r' in the next guess but not in the same position, and avoid 's', 't', 'e'.  } \\
\texttt{guess: rater} \\
            }
        }
    }
    & & \\ \\

    \theutterance \stepcounter{utterance}  
    & & & \multicolumn{2}{p{0.3\linewidth}}{
        \cellcolor[rgb]{0.9,0.9,0.9}{
            \makecell[{{p{\linewidth}}}]{
                \texttt{\tiny{[GM$|$GM]}}
                \texttt{attempts = 2} \\
\texttt{target\_word = crank} \\
\texttt{guess = rater} \\
\texttt{guess\_feedback = r<yellow> a<yellow> t<red> e<red> r<red>} \\
            }
        }
    }
    & & \\ \\

    \theutterance \stepcounter{utterance}  
    & & & \multicolumn{4}{p{0.6\linewidth}}{
        \cellcolor[rgb]{0.9,0.9,0.9}{
            \makecell[{{p{\linewidth}}}]{
                \texttt{\tiny{[P1$\langle$GM]}}
                \texttt{guess\_feedback: r<yellow> a<yellow> t<red> e<red> r<red>} \\
\\ 
\texttt{Provide your response only in this format.} \\
\texttt{explanation: details} \\
\texttt{guess: word} \\
            }
        }
    }
     \\ \\

    \theutterance \stepcounter{utterance}  
    & \multicolumn{4}{p{0.45\linewidth}}{
        \cellcolor[rgb]{0.95,0.95,0.95}{
            \makecell[{{p{\linewidth}}}]{
                \texttt{\tiny{[P1$|internal]}}
                \texttt{{[}attempt 1 – valid{]}} \\
\texttt{explanation: The first guess revealed that 'a' is in the correct position (green), and 'r' is present but in the wrong position (yellow). The letters 's', 't', and 'e' are incorrect (red). From the second guess, 'r' is still yellow, meaning it's in the word but not in the fourth position, and 'a' remains yellow, indicating it's in the word but not in the second position. 't', 'e', and 'r' in the second guess are red, meaning they are not in the word. However, 'r' is confirmed to be in the word from the first guess, so the red 'r' in the second guess must be a false negative due to position. The correct positions are: 'a' is in position 3, and 'r' is in the word but not in position 4. The second guess also shows 'a' is not in position 2, so 'a' must be in position 3. The word has 'a' in position 3 and 'r' in position 1 or 5. Since 'r' is yellow in the first guess, it's not in position 4. The second guess shows 'r' in position 1, but it's yellow, so 'r' is in the word but not in position 1. Therefore, 'r' must be in position 5. The word has '} \\
            }
        }
    }
    & & \\ \\

    \theutterance \stepcounter{utterance}  
    & \multicolumn{4}{p{0.6\linewidth}}{
        \cellcolor[rgb]{0.561,0.737,0.561}{
            \makecell[{{p{\linewidth}}}]{
                \texttt{\tiny{[P1$\rangle$GM]}}
                \texttt{explanation: The first guess revealed that 'a' is in the correct position (green), and 'r' is present but in the wrong position (yellow). The letters 's', 't', and 'e' are incorrect (red). From the second guess, 'r' is still yellow, meaning it's in the word but not in the fourth position, and 'a' remains yellow, indicating it's in the word but not in the second position. 't', 'e', and 'r' in the second guess are red, meaning they are not in the word. However, 'r' is confirmed to be in the word from the first guess, so the red 'r' in the second guess must be a false negative due to position. The correct positions are: 'a' is in position 3, and 'r' is in the word but not in position 4. The second guess also shows 'a' is not in position 2, so 'a' must be in position 3. The word has 'a' in position 3 and 'r' in position 1 or 5. Since 'r' is yellow in the first guess, it's not in position 4. The second guess shows 'r' in position 1, but it's yellow, so 'r' is in the word but not in position 1. Therefore, 'r' must be in position 5. The word has '} \\
            }
        }
    }
    & & \\ \\

    \theutterance \stepcounter{utterance}  
    & & & \multicolumn{2}{p{0.3\linewidth}}{
        \cellcolor[rgb]{0.9,0.9,0.9}{
            \makecell[{{p{\linewidth}}}]{
                \texttt{\tiny{[GM$|$GM]}}
                \texttt{Error: The response should contain the 'guess:' keyword exactly once.} \\
            }
        }
    }
    & & \\ \\

    \theutterance \stepcounter{utterance}  
    & & & \multicolumn{2}{p{0.3\linewidth}}{
        \cellcolor[rgb]{0.9,0.9,0.9}{
            \makecell[{{p{\linewidth}}}]{
                \texttt{\tiny{[GM$|$GM]}}
                \texttt{game\_result = ABORT} \\
            }
        }
    }
    & & \\ \\

\end{supertabular}
}

\end{document}
